% LSPV-1 REPORT::LEAF2::2026-09-30T19:40:16-07:00
\documentclass[11pt]{article}
\usepackage[margin=0.9in]{geometry}
\usepackage{amsmath,amssymb}
\begin{document}
\section*{LSPV-1: RMV leaf reduction, extractor counting, and univariate rank battery}

\textbf{STATUS: OPEN.} P versus NP remains open. This run narrows the SAME\_REALIZATION product-span question but does not close it.

\paragraph{PROVED/EXACT: every ordinary axis-parallel univariate restriction occurs at an RMV leaf.}
Definition 22 recursively enumerates every MV line in both half-directions and every fixing value. Since the dimension $r$ is a power of two, fix a coordinate $c$ and values for all other coordinates. At each recursive level choose the direction containing $c$ and fix the complementary half to the prescribed values. At the final $r=2$ call choose the final one-dimensional direction. Hence the resulting ordinary axis-parallel restriction
\[
t\longmapsto \widehat g(a_1,\ldots,a_{c-1},t,a_{c+1},\ldots,a_r)
\]
is one of the univariate polynomial oracles fed to the local extractor and all of its seed outputs are contained in $RMV(\widehat g)$.

\paragraph{PROVED/EXACT: any nonzero canonical PCP proof has a nonconstant leaf.}
The vMS canonical LDE equals the proof bit on represented grid points in $[h]^r$ and is zero on unused grid points. Theorem 24 chooses $r$ so the grid strictly exceeds the proof length. Thus if a perfectly accepted proof $y$ contains a $1$, the grid values contain both $0$ and $1$. Connect such two grid points by changing one coordinate at a time. Along that coordinate path some axis-parallel grid line changes value, and its univariate restriction is a degree-at-most $h-1$ polynomial whose $h$ grid values form a nonconstant Boolean word. By the previous paragraph, that leaf occurs inside the same RMV realization. The all-zero accepted-proof case remains a genuine degenerate control.

\paragraph{PROVED/EXACT: extractor bad-rank counting lemma.}
Let $E$ be a $(\kappa,\varepsilon)$ local $\mathcal C$-extractor with $m$ output bits and $\varepsilon<1/4$. For $f\in\mathcal C$ write
\[
H_f=\{E^f(y):y\},\qquad M=1+m+\binom m2,
\]
and let $\rho(f)$ be the rank of degree-$\le2$ Boolean features on $H_f$. If $\rho(f)<M$, some nonzero Boolean quadratic $Q$ vanishes on $H_f$. Every nonzero Boolean quadratic has relative weight at least $1/4$, so its zero set $D_Q$ has density at most $3/4$. Therefore
\[
\Pr_y[E^f(y)\in D_Q]=1> |D_Q|/2^m+\varepsilon.
\]
Shaltiel--Umans Proposition 3.8 implies that for each fixed $Q$ at most $2^\kappa$ source objects $f$ can satisfy this. Union bounding over the fewer than $2^M$ nonzero quadratics gives
\[
\boxed{|\{f\in\mathcal C:\rho(f)<M\}|<2^{\kappa+M}.}
\]
Consequently any subset of the source class containing more than $2^{\kappa+M}$ distinct objects contains a full-quadratic-rank element. This is structural only; it is not a runtime lower bound.

\paragraph{Important limitation of the counting lemma.}
For the vMS local extractor $\kappa=\Delta^5$. At univariate RMV leaves the entire degree-$\le\Delta$ polynomial class has only $p^{\Delta+1}$ elements, and with $p\asymp\Delta^{100}$ one has $\log_2|\mathcal C|=O(\Delta\log\Delta)\ll\Delta^5$. Thus the extractor guarantee is vacuous at the leaf for this counting purpose. The lemma becomes informative only at sufficiently high-dimensional RM nodes. It does not prove the leaf conjecture below.

\paragraph{COMPUTATIONAL: nonconstant-leaf full-rank conjecture survives a larger battery.}
For the declared SU05 Eq.~(7) univariate orbit plus standard binary-label Hadamard realization, new exact finite tests interpolate every nonconstant Boolean word on $h$ grid points and compute the exact quadratic-feature rank over all nonzero orbit seeds and all Hadamard coordinates. The following exhaustive cases all had rank exactly $M=1+m+\binom m2$ for every nonconstant Boolean word:
\[
(p,h,m)=(17,4,4),(29,5,5),(37,6,6),(43,7,7),(53,8,8).
\]
For $(p,h,m)=(31,6,8)$, all $62$ nonconstant words had full rank for each of all eight primitive roots modulo $31$. Selected adversarial patterns (singleton, adjacent pair, alternating, half-block, and a modular pattern) also had full rank at $(257,16,16)$, $(257,16,32)$, $(257,32,32)$, and $(257,64,64)$, the last having $M=2081$. No deficient case was found. This remains \textbf{COMPUTATIONAL}, not a theorem.

\paragraph{FALSIFIED/RETIRED simplification: one Hadamard bit need not suffice.}
A tempting reduction is to fix one Hadamard basis coordinate and analyze one binary cyclic sequence. It fails as a generic proof strategy: at $(p,h,m)=(31,6,8)$ with $M=37$, the five individual label-bit restrictions of tested patterns had ranks only in the range $23$--$30$, although the union over Hadamard coordinates had full rank $37$. Cross-bit quadratic information is genuinely useful.

\paragraph{Literature/model interpretation.}
SU09 Proposition 3.8 is already an image-size theorem for a local extractor and directly yields the counting lemma above. The Schur-square literature studies the same rank object: random codes usually acquire maximal square dimension, while algebraic or cyclic codes may have constrained squares. Recent cyclic/constacyclic work computes Schur powers by Fourier/additive-combinatorial descriptions. Carlini--Catalisano--Geramita study exactly the degree-two Hilbert function of generic unions of linear spaces, the algebraic-geometric analogue of the Hadamard-block arrangement $H=\bigcup_v S_v$. Chizhov analyzes square-code dimension through the rank of matrices of quadratic-form evaluations, which is especially close to our $\Phi_2(S_v)$ formulation.

The BGH+SV robust PCPP source does not currently supply the hoped-for low-rank theorem. For general circuit satisfiability its proof machinery explicitly contains oracles evaluating linear and quadratic functions of the computation tableau, whereas it becomes linear only for linear circuits. This is evidence that perfect PCP validity alone should not be assumed to align the $S_v$ quadratically.

\paragraph{Strict status and next target.}
PROVED/EXACT: RMV leaf coverage; nonzero-proof $\Rightarrow$ nonconstant Boolean-grid leaf; extractor bad-rank counting lemma. COMPUTATIONAL: every tested nonconstant SU05/Hadamard leaf has maximal quadratic rank. FALSIFIED/RETIRED: single-Hadamard-bit reduction as a general route. OPEN: the all-parameter nonconstant-leaf theorem for the exact permitted binary encoder/local extractor; the all-zero accepted-proof control; arbitrary-algorithm transfer; acquisition; relativization/algebrization; final diagonal.

\textbf{Single next target:} write the quadratic-relation equations for the sliding matrices $B_v$ and prove or break: every nonconstant Boolean-grid interpolant in the declared SU05/Hadamard leaf model has $\rho=M$. The closest literature attack is via cyclic-code Schur squares / DFT and quadratic-form evaluation matrices, not generic ANF expansion.

\paragraph{Artifacts.}
\texttt{experiments/product\_span/su05\_leaf\_line\_rank.py};
\texttt{results/raw/product\_span/2026-09-30\_LSPV1\_su05\_leaf\_line\_rank.json}.

\begin{thebibliography}{9}
\bibitem{vMS23} D. van Melkebeek and N. Mocelin Sdroievski, ``Instance-Wise Hardness Versus Randomness Tradeoffs for Arthur-Merlin Protocols,'' CCC 2023, especially Lemma 18 and Definitions 20--22.
\bibitem{SU09} R. Shaltiel and C. Umans, ``Low-End Uniform Hardness Versus Randomness Tradeoffs for AM,'' SIAM J. Comput. 39 (2009). See the local extractor definition, Proposition 3.8, and RMV construction.
\bibitem{SU05} R. Shaltiel and C. Umans, ``Simple Extractors for All Min-Entropies and a New Pseudorandom Generator,'' J. ACM 52(2):172--216, 2005.
\bibitem{BGHSV} E. Ben-Sasson, O. Goldreich, P. Harsha, M. Sudan, S. Vadhan, ``Robust PCPs of Proximity, Shorter PCPs, and Applications to Coding,'' SIAM J. Comput. 36(4):889--974, 2006.
\bibitem{CCG} E. Carlini, M. V. Catalisano, A. V. Geramita, ``Subspace arrangements, configurations of linear spaces and the quadrics containing them,'' J. Algebra 362:70--83, 2012.
\bibitem{Chizhov} I. V. Chizhov, ``Hadamard square of linear codes and the generalized minimal distance of Reed--Muller code of order 2,'' Diskr. Mat. 35(1):128--152, 2023.
\end{thebibliography}
\end{document}
% END_OF_VERIFIED_REPORT::LSPV-1::LEAF2
